\documentclass[11pt,a4paper]{article}
\usepackage[margin=25mm]{geometry}
\usepackage{amsmath,amssymb,amsthm,booktabs}
\newtheorem{lemma}{Lemma}\newtheorem{theorem}{Theorem}
\title{Two pair counts with an element of order twelve}
\author{Human Compression project}\date{}
\begin{document}\maketitle
Let $G=O'N$, and write $x_J(p,q)$ for the number of ordered pairs
$(r,s)$ in $J$ with $|r|=p$, $|s|=q$ and $|rs|=2$. Let $e_J(p,q)$
count those pairs which generate $J$. We prove
\[
       e_G(7,12)/(2|G|)=815,\qquad e_G(8,12)/(2|G|)=2370.       \tag{1}
\]
These are the fixed Connor--Leemans Table~1 entries. The complete subgroup
roster, the ATLAS centralizer and named maximal subgroup data, and the
previously acquired character evaluations remain explicit borrowed
premises. Row numbers below refer to that roster.

\begin{lemma}
For an element $s$ of order twelve, exactly one maximal subgroup in
row 7 contains $s$. If $(r,s)$ has type $(7,12)$ or $(8,12)$ and generates
a proper subgroup not contained in this maximal subgroup, then it
has a unique maximal overgroup in $G$.
\end{lemma}
\begin{proof}
The ATLAS centralizer has order 36. Its central subgroup
$\langle s^3\rangle$ of order four is a Sylow 2-subgroup. A central
Sylow subgroup is normal, hence the unique Sylow subgroup of its
prime. Thus $C_G(s)$ has just one involution, $t=s^6$.
The row 7 maximal subgroups are the involution centralizers, so
its unique member containing $s$ is $C(s)=C_G(t)$.

The cyclic order seven classes are rows 569 and 570, the cyclic order
eight classes are 559 and 560, and the cyclic order twelve class is
541. For each of the two types, intersect the relevant complete ancestor
lists, remove $G$ and every class which occurs below row 7. The
remaining rows are respectively
\[
\begin{array}{c|l}
(7,12)&2,3,21,22,29,68\\
(8,12)&2,3,11,12,18,19,21,26,38,40,56,98.
\end{array}                                                     \tag{2}
\]
In the first list, rows 2, 3 and 21 are maximal. Row 29 has only
row 22 as an immediate overgroup class, and row 68 has only row 21.
The direct incidences between class lengths and cover multiplicities
show that each of these parents is unique as an actual subgroup.
The retained literal flag normalizer $B$ in row 22 has $x_B(7,12)=0$,
as verified below. This excludes both rows 22 and 29. Every remaining
possibility has a unique maximal overgroup.

In the second list, rows 2, 3, 11, 12 and 21 are maximal. The complete
immediate-parent lists of the other rows are
\[
18\longrightarrow11,\quad19\longrightarrow11,\quad
26\longrightarrow12,\quad38,40\longrightarrow26,\quad
56\longrightarrow21,\quad98\longrightarrow56.
\]
For every displayed edge, the number of parent conjugates times the
number of child subgroups in a fixed parent equals the number of
child conjugates. Double counting incidences therefore proves one
actual parent at each step. Iterating these uniqueness statements
proves the required unique maximal overgroup. Arbitrary cover paths
are not treated as distinct containments.

Finally, a subgroup class below row 7 has an actual row 7 overgroup
by conjugacy. If it contains $s$, that overgroup must be $C(s)$.
Thus the roster exclusion is precisely the condition in the statement.
\end{proof}

We apply the accepted distinguished maximal counting lemma. A pair
outside $C(s)$ occurs just once in the maximal sum. A pair inside
$C(s)$ is overcounted once for each other maximal overgroup containing
it. The corrections count exactly the pairs in those other maximal
subgroups for which $r$ commutes with $s^6$.

\begin{lemma}
For $p\in\{7,8\}$, the correction to the maximal sum is
\[
            2\cdot342868680\,x_H(p,12)
                     +300010095\,x_T(p,12),                   \tag{3}
\]
where $H$ is the accepted block group of order 1344 and $T$ is the
accepted kernel-involution centralizer of order 1536. Their values are
\[
\begin{array}{c|rr}
p&7&8\\\hline
x_H(p,12)&672&1344\\
x_T(p,12)&0&6144.
\end{array}                                                     \tag{4}
\]
\end{lemma}
\begin{proof}
The only maximal classes with nonzero $x_M(7,12)$ are rows 2, 3, 7
and 21; for $(8,12)$ there is the additional row 12. In either
projective maximal $M\cong L_3(7):2$, the involution $s^6$ belongs
to the socle, since its image in the quotient of order two is trivial.
The accepted block centralizer proof gives one class of inner
involutions, with centralizer $H$. There are $3753792/1344=2793$
such involutions in each $M$, and 122760 actual maximal subgroups in
each of these two classes. Their product is $342868680$.
Summing over the involutions counts precisely the required pairs:
an involution commuting with the second entry must equal $s^6$.

For row 21, reuse the accepted 10752-element model with normal kernel
$C_4^3$. Its seven nonidentity kernel involutions form a single orbit,
with centralizer $T$. The literal reader checks all 1792 elements of
order twelve and verifies that their sixth powers lie in this orbit.
The number of actual row 21 maximal subgroups is 42858585, so this
contribution is $42858585\cdot7\,x_T(p,12)=300010095\,x_T(p,12)$.
Since $7$ does not divide $1536$, its first value in (4) is zero.

It remains to exclude a correction from row 12 for $(8,12)$.
The complete roster has no subgroup class which both contains
elements of orders eight and twelve and occurs below both rows
7 and 12. This is checked by intersecting the two cyclic ancestor
sets with the two descendant sets; the result is empty. If a pair
in a row 12 maximal had $r$ commuting with $s^6$, its generated
subgroup would be contained in that maximal and in $C(s)$, contradicting
this empty intersection. Thus row 12 contributes no correction.
\end{proof}

\paragraph{The finite certificate.}
All groups are reused from accepted literal permutation models.
The independent reader reconstructs each complete group from its
retained generators, partitions its order twelve elements into
conjugacy classes, and checks every partner of every representative.
It also records whether a partner commutes with the sixth power.
In $B$, the only partner orders are six and twelve, proving the
exclusion used in the first lemma. The calculations in $H$ and $T$
give (4). The reader reconstructs the 10752-element model and checks
the entire sixth-power assertion. It uses pointwise permutation
composition independently of the acquisition program's implementation.
The complete perfect-group catalogue and the original identification
of that model remain the explicit premises of the accepted order
twelve exclusion; they are not reproved here.

\begin{theorem}
For $p\in\{7,8\}$, with $M$ ranging over representatives of all thirteen
maximal subgroup classes,
\[
 e_G(p,12)=x_G(p,12)-\sum_M[G:M]x_M(p,12)
       +2\cdot342868680\,x_H(p,12)+300010095\,x_T(p,12).
\]
The ordered automorphism-orbit counts are those in (1).
\end{theorem}
\begin{proof}
The two lemmas and the distinguished maximal identity give the formula.
Its character-table evaluations and the values (4) are exact integers.
Division by $|\operatorname{Aut}(G)|=2|G|$ is valid because an
automorphism fixing a generating pair fixes $G$ pointwise. The
certificate evaluates the formula and compares both answers with
the fixed printed entries. No extra quotient by duality is taken.
\end{proof}

The numerical assertions are due to Connor and Leemans. The present
contribution is their derivation by the shared counting lemma, exact
overgroup joins and the small literal centralizer calculation.
These are two additional numerical entries. They supply no new chirality
statement, source completion, open-question answer or Lean certificate.
\end{document}
